\section{Modelo de datos y evaluación de factibilidad}
\label{sec:modelo}

\subsection{Instancia}

Una instancia (\cod{Instance}) se lee en el formato de Solomon~\cite{solomon1987}, que también
usan las instancias de Gehring y Homberger~\cite{gehring1999}. Contiene el depósito (nodo $0$) y
$n$ clientes ($N=n+1$ nodos). Cada nodo $i$ tiene coordenadas $(x_i,y_i)$, demanda $\mu_i$,
ventana de tiempo $[a_i,b_i]$ (\cod{ready\_time}, \cod{due\_date}) y tiempo de servicio $s_i$.
Todos los vehículos tienen capacidad $Q$. El archivo declara además un tamaño de flota, que el
solver lee pero \emph{no impone} como restricción: el número de vehículos es un objetivo, no
un límite. La flota con la que trabaja cada etapa de la búsqueda la fija el propio algoritmo
(Sección~\ref{sec:fases}).

Al construir la instancia se precalculan dos matrices $N\times N$ que se consultan en tiempo
constante durante toda la búsqueda:
\begin{align}
d_{ij} &= \sqrt{(x_i-x_j)^2+(y_i-y_j)^2}, \label{eq:dist}\\
\alc(i,j) &= \bigl[\, i\neq j \ \wedge\ a_i+s_i+d_{ij}\le b_j \,\bigr]. \label{eq:alc}
\end{align}
La distancia~\eqref{eq:dist} es euclidiana en doble precisión, sin redondeo ni truncamiento, y
coincide con el tiempo de viaje. La matriz de alcanzabilidad~\eqref{eq:alc} indica si es
posible visitar $j$ inmediatamente después de $i$ en el mejor caso (saliendo de $i$ lo antes
posible); si $\alc(i,j)$ es falso, el arco $(i,j)$ no puede aparecer en ninguna ruta
factible. Es una condición necesaria, no suficiente, y se usa como filtro previo barato.
Durante el mismo precálculo se guarda la mayor distancia entre dos nodos,
$d_{\max}=\max_{i,j}d_{ij}$ (\cod{max\_dist}), que fija la amplitud del ruido de los
operadores de reparación (Sección~\ref{sec:ruido}).

La instancia guarda además una tercera estructura, \cod{shaw\_order}, que se construye de
forma perezosa la primera vez que se necesita y se describe en la
Sección~\ref{sec:shaw}.

\subsection{Ruta}

Una ruta (\cod{Route}) es una secuencia $\pi=(\pi_0,\pi_1,\dots,\pi_{L-1})$ con
$\pi_0=\pi_{L-1}=0$. Una ruta con $L=2$ está vacía y no cuenta como vehículo usado. Junto a
la secuencia se almacenan su carga, su distancia y tres vectores de longitud $L$:
\begin{itemize}
\item $A_k$: instante de llegada a $\pi_k$ (\cod{arrival\_times});
\item $W_k=\max(0,\,a_{\pi_k}-A_k)$: espera en $\pi_k$ hasta que abre su ventana (\cod{wait\_times});
\item $F_k$: holgura temporal hacia adelante (\cod{time\_slacks}), es decir, cuánto puede
retrasarse el inicio del servicio en $\pi_k$ sin violar la ventana de $\pi_k$ ni la de ningún
nodo posterior de la ruta.
\end{itemize}
El inicio del servicio en $\pi_k$ es $\max(A_k,a_{\pi_k})$. La holgura sigue la recurrencia
clásica de Savelsbergh~\cite{savelsbergh1992}:
\begin{equation}
F_{L-1}=b_{0}-\max(A_{L-1},a_{0}),\qquad
F_k=\min\bigl(b_{\pi_k}-\max(A_k,a_{\pi_k}),\; F_{k+1}+W_{k+1}\bigr).
\label{eq:slack}
\end{equation}
El segundo término expresa que un retraso en $\pi_k$ se transmite al nodo siguiente, donde
primero consume la espera $W_{k+1}$ (tiempo que de todos modos se iba a perder) y solo el
excedente consume su holgura $F_{k+1}$.

El Algoritmo~\ref{alg:recalcular} recalcula todos estos valores con una pasada hacia adelante
y otra hacia atrás, en tiempo $\bigO(L)$. Se invoca cada vez que una ruta cambia.

\begin{algorithm}[H]
\caption{\fn{RecalcularRuta} (\cod{Route::recalculate})}
\label{alg:recalcular}
\Entrada{ruta $\pi$ de longitud $L$, instancia}
\Salida{carga, distancia y vectores $A$, $W$, $F$ de la ruta}
$\textit{carga}\gets 0$;\quad $\textit{dist}\gets 0$;\quad $A_0\gets 0$;\quad $W_0\gets 0$\;
\Para{$k\gets 0$ \Hasta $L-2$}{
  $u\gets\pi_k$;\quad $v\gets\pi_{k+1}$\;
  $\textit{dist}\gets\textit{dist}+d_{uv}$\;
  \lSi{$u\neq 0$}{$\textit{carga}\gets\textit{carga}+\mu_u$}
  $A_{k+1}\gets\max(A_k,a_u)+s_u+d_{uv}$ \Com*[r]{inicio de servicio en $u$ + servicio + viaje}
  $W_{k+1}\gets\max(0,\,a_v-A_{k+1})$\;
}
$F_{L-1}\gets b_{0}-\max(A_{L-1},a_{0})$\;
\Para{$k\gets L-2$ \Hasta $0$}{
  $F_k\gets\min\bigl(b_{\pi_k}-\max(A_k,a_{\pi_k}),\;F_{k+1}+W_{k+1}\bigr)$\;
}
\end{algorithm}

\subsection{Evaluación de una inserción en tiempo constante}
\label{sec:evalins}

La operación más frecuente de todo el solver es preguntar si un cliente $u$ puede insertarse
entre $\pi_k$ y $\pi_{k+1}$ y cuánto aumenta la distancia. Con los vectores $A$, $W$ y $F$ la
respuesta se obtiene en $\bigO(1)$, sin recorrer el resto de la ruta
(Algoritmo~\ref{alg:evalins}).

\begin{algorithm}[H]
\caption{\fn{EvaluarInserción} (\cod{evalInsertion})}
\label{alg:evalins}
\Entrada{cliente $u$, ruta $\pi$ con $A,W,F$ vigentes, posición $k$}
\Salida{(\textit{factible}, $\Delta$)}
$i\gets\pi_k$;\quad $j\gets\pi_{k+1}$\;
\lSi{$\neg\alc(i,u)$ \O $\neg\alc(u,j)$}{\Devolver (falso, --)}
$A_u\gets\max(A_k,a_i)+s_i+d_{iu}$ \Com*[r]{llegada a $u$}
\lSi{$A_u>b_u$}{\Devolver (falso, --)}
$A'_j\gets\max(A_u,a_u)+s_u+d_{uj}$ \Com*[r]{nueva llegada a $j$}
$\textit{retraso}\gets\max(0,\,A'_j-A_{k+1})$\;
\lSi{$\textit{retraso}>W_{k+1}+F_{k+1}$}{\Devolver (falso, --)}
$\Delta\gets d_{iu}+d_{uj}-d_{ij}$\;
\Devolver (verdadero, $\Delta$)\;
\end{algorithm}

La corrección de la prueba se sigue de la definición de $F$: el retraso en la llegada a $j$
es absorbido primero por la espera $W_{k+1}$, y el inicio del servicio en $j$ se desplaza en
$\max(0,\textit{retraso}-W_{k+1})$, que por~\eqref{eq:slack} es admisible para $j$ y para todos
los nodos posteriores si y solo si no supera $F_{k+1}$. La restricción de capacidad no depende
de la posición, por lo que se verifica una sola vez por ruta
($\textit{carga}+\mu_u\le Q$) antes de recorrer sus posiciones.

Evaluar todas las posiciones de una ruta cuesta entonces $\bigO(L)$ en lugar de
$\bigO(L^2)$, y evaluar un cliente contra toda la solución cuesta $\bigO(n+m)$. El precio es
que cada modificación de una ruta exige un recálculo $\bigO(L)$
(Algoritmo~\ref{alg:recalcular}), que es despreciable frente al número de evaluaciones que
habilita.

\subsection{Solución y función de costo}

Una solución (\cod{Solution}) $x$ contiene el vector de rutas $\Rset(x)$, la lista $\Uns(x)$ de
clientes sin asignar (\cod{unassigned}) y dos métricas agregadas que mantiene
\cod{updateMetrics}: el número de vehículos $K(x)$, que cuenta solo las rutas no vacías, y la
distancia total $D(x)$ de esas rutas. Una solución es \emph{completa} si $\Uns(x)=\emptyset$.

Todas las rutas son factibles en todo momento: los operadores de destrucción solo quitan
clientes (lo que nunca rompe la factibilidad cuando se cumple la desigualdad triangular) y los
de reparación solo ejecutan inserciones aceptadas por el Algoritmo~\ref{alg:evalins}. La
búsqueda nunca visita soluciones con violaciones de capacidad o de ventanas; la única forma de
infactibilidad posible es que queden clientes sin asignar, y la función de costo la penaliza
de forma explícita.

El objetivo jerárquico habitual del VRPTW (primero minimizar vehículos, después distancia) se
escalariza con un peso fijo por vehículo, al que se antepone una penalización por cada
cliente sin asignar:
\begin{equation}
f(x)=10^{6}\cdot|\Uns(x)|+10^{4}\cdot K(x)+D(x).
\label{eq:costo}
\end{equation}
Mientras los pesos dominen a los términos siguientes (en particular, mientras $D(x)<10^4$ por
cada vehículo de diferencia), comparar $f$ equivale a la comparación lexicográfica
$(|\Uns|,K,D)$: primero colocar a todos los clientes, después usar menos vehículos y por
último recorrer menos distancia. Gracias al primer término, $f$ está definida también para
soluciones incompletas, que por tanto participan en la búsqueda como cualquier otra; esto es
lo que permite la fase de minimización de rutas de la Sección~\ref{sec:fases}. La función $f$
es la que usan el criterio de aceptación y la detección de mejoras; las ganancias que
alimentan el aprendizaje se miden con el mismo orden lexicográfico, pero en términos
relativos (Sección~\ref{sec:ganancia}).

\subsection{Solución inicial}

La solución inicial se construye con una heurística secuencial del vecino más cercano
(Algoritmo~\ref{alg:inicial}): cada ruta parte del depósito y agrega repetidamente el cliente
no visitado más cercano al último nodo que sea alcanzable, respete la capacidad, llegue dentro
de su ventana y permita regresar al depósito antes de $b_0$. Cuando ningún cliente cumple
estas condiciones la ruta se cierra y se abre otra. Es determinista, por lo que ambos
algoritmos y todas las semillas parten de la misma solución; su calidad es deliberadamente
modesta y toda la mejora posterior corresponde al ALNS.

\begin{algorithm}[H]
\caption{\fn{SoluciónInicial} (\cod{Solution::generateInitialSolution})}
\label{alg:inicial}
\Entrada{instancia}
\Salida{solución $x$ con rutas factibles}
\Mientras{queden clientes sin visitar}{
  $\pi\gets(0,0)$;\quad $c\gets 0$;\quad $t\gets 0$;\quad $\textit{carga}\gets 0$\;
  \Mientras{verdadero}{
    $\mathcal{C}\gets\{\,i\text{ no visitado}:\ \alc(c,i),\ \textit{carga}+\mu_i\le Q,\ t_i\le b_i,\ t_i+s_i+d_{i0}\le b_0\,\}$\;
    \hspace{2em}donde $t_i=\max(t+s_c+d_{ci},\,a_i)$\;
    \uSi{$\mathcal{C}\neq\emptyset$}{
      $i^\star\gets\arg\min_{i\in\mathcal{C}} d_{ci}$\;
      insertar $i^\star$ al final de $\pi$; marcarlo visitado\;
      $\textit{carga}\gets\textit{carga}+\mu_{i^\star}$;\quad $t\gets t_{i^\star}$;\quad $c\gets i^\star$\;
    }
    \SiNo{
      \lSi{$c=0$}{enviar a $\Uns$ todos los clientes no visitados \Com*[f]{no caben ni solos}}
      \Salir\;
    }
  }
  \fn{RecalcularRuta}($\pi$);\quad agregar $\pi$ a $\Rset(x)$\;
}
actualizar $K(x)$ y $D(x)$\;
\end{algorithm}
